\documentclass{beamer}
\usepackage[utf8]{inputenc}
\usepackage[brazilian]{babel}
\usepackage{booktabs}
\usepackage{xcolor}

% Tema CERISE — cores da identidade visual do projeto (ver
% .claude/skills/atualizar-slide-migracao-langchain/references/tema_cerise.sty).
% Colado direto aqui para o arquivo funcionar sozinho no Overleaf sem precisar
% subir um segundo arquivo .sty.
\definecolor{cerisemagenta}{HTML}{A6186B}
\definecolor{ceriserosa}{HTML}{E6186D}
\definecolor{cerisecoral}{HTML}{F5642D}
\definecolor{ceriseroxo}{HTML}{5B1A6B}

\setbeamercolor{title}{fg=cerisemagenta}
\setbeamercolor{frametitle}{fg=white,bg=cerisemagenta}
\setbeamercolor{structure}{fg=cerisemagenta}
\setbeamercolor{block title}{fg=white,bg=ceriseroxo}
\setbeamercolor{block body}{bg=cerisemagenta!5}
\setbeamercolor{alerted text}{fg=cerisecoral}
\setbeamercolor{item}{fg=cerisemagenta}

\newenvironment{numerobloco}[1]{\begin{block}{#1}}{\end{block}}
\newenvironment{decisaobloco}[1]{\begin{alertblock}{#1}}{\end{alertblock}}

\title{Migração do Harbor para LangChain/LangGraph}
\subtitle{Progresso, comparações e decisões}
\author{Yan Santos Leite}
\date{\today}

\begin{document}

\frame{\titlepage}

\begin{frame}{Por que migrar}
  \begin{itemize}
    \item O roteador (\texttt{dashboard/roteador.py}) e o self-repair do NL-to-SQL
      (\texttt{DBA-Agent}) usam \textbf{LangGraph} para branching condicional e um padrão
      de \emph{critic node} (CRAG/Self-RAG).
    \item Ter o RAG já em LangChain evitou uma segunda migração quando esses módulos
      passaram a consumir o retriever dentro de um grafo.
    \item Cada módulo só migra depois de medido — nunca ``porque LangChain é melhor'' sem
      número que sustente.
  \end{itemize}
\end{frame}

\begin{frame}{Ordem de migração (por risco crescente) — estado em 2026-09-09}
  \begin{enumerate}
    \item \textbf{RAG híbrido} — \alert{migrado} (2026-09-08)
    \item \textbf{Self-repair / DBA-Agent} — \alert{migrado} (2026-09-09), sem regressão
    \item \textbf{Roteador + answerability gates} — \alert{migrado} (2026-09-09), 0
      divergências em 67 perguntas do golden set
    \item Diagnóstico em camadas (regra $\to$ Isolation Forest $\to$ LLM) — não iniciado,
      baixo retorno esperado
    \item NL-to-SQL (geração em si) — não iniciado, maior risco (validação determinística
      própria precisa ser preservada)
  \end{enumerate}
\end{frame}

\begin{frame}{RAG híbrido: TF-IDF vs. BM25}
  \small
  \begin{numerobloco}{O que medimos (golden set, 49 perguntas RAG, 2026-08-22)}
    \begin{center}
    \begin{tabular}{lccc}
      \toprule
      \textbf{Métrica} & \textbf{TF-IDF (produção)} & \textbf{BM25} \\
      \midrule
      Recall@5 (hybrid+rerank) & 98\%   & 98\%   \\
      Precision@5 (hybrid+rerank) & 44\% & 43\%  \\
      MRR (hybrid+rerank)      & 0,866  & 0,866  \\
      Recall (lexical isolado) & 98\%   & 94\%   \\
      \bottomrule
    \end{tabular}
    \end{center}
  \end{numerobloco}

  \pause
  \textbf{O que isso significou:} empate estatístico no pipeline completo — o Cross-Encoder
  absorve a diferença que aparece no lexical isolado (BM25 perde: 94\% vs 98\%). BM25 tem
  $\sim$15\% menos latência, mas isso não compensa qualidade pior no isolado.

  \pause
  \begin{decisaobloco}{O que decidimos}
    Mantivemos TF-IDF como \emph{default} — o resultado não justificou a troca. BM25
    continua disponível via \texttt{lexico="bm25"} se o corpus crescer bastante.
  \end{decisaobloco}
\end{frame}

\begin{frame}{RAG híbrido: Python puro vs. LangChain}
  \small
  \begin{numerobloco}{O que medimos (golden set, 49 perguntas de rota RAG, 2026-09-08)}
    \begin{center}
    \begin{tabular}{lccc}
      \toprule
      \textbf{Métrica} & \textbf{Prod. (TF-IDF)} & \textbf{LC fiel} & \textbf{LC BM25+RRF} \\
      \midrule
      Recall@5    & 98,0\%  & 98,0\%  & 98,0\%  \\
      Precision@5 & 38,4\%  & 38,4\%  & \alert{49,8\%} \\
      MRR         & 0,898   & 0,898   & 0,897   \\
      \bottomrule
    \end{tabular}
    \end{center}
  \end{numerobloco}

  \pause
  \textbf{O que isso significou:} o ganho de qualidade (+11,4pp de Precision@5) vem do
  \alert{algoritmo} BM25+RRF, não do framework LangChain — a versão ``fiel'' (mesmo algoritmo,
  framework trocado) empata exatamente com a produção em Recall@5/MRR.

  \pause
  \begin{decisaobloco}{O que decidimos}
    Migramos mesmo assim, por consistência com o roteador e o self-repair (ambos LangGraph
    logo em seguida) — não pelo ganho isolado de retrieval. Usamos a implementação BM25+RRF
    (framework + o ganho real juntos), não a versão ``fiel''.
  \end{decisaobloco}
\end{frame}

\begin{frame}{Self-repair NL-to-SQL: Python puro vs. LangGraph}
  \small
  \begin{numerobloco}{O que medimos (harness completo, 67 perguntas, 2026-09-09)}
    \begin{center}
    \begin{tabular}{lcc}
      \toprule
      \textbf{Métrica} & \textbf{Antes (baseline)} & \textbf{Depois (LangGraph)} \\
      \midrule
      Roteamento correto & 56/67  & 55/67$^*$ \\
      Faithfulness médio & 64\%   & 64\%   \\
      Alucinações         & 4      & 3      \\
      \bottomrule
    \end{tabular}
    \end{center}
    \footnotesize $^*$ pergunta de fronteira (\texttt{cnc-cam3}) oscila mesmo no roteador
    \emph{original} — instabilidade do LLM, não regressão.
  \end{numerobloco}

  \pause
  \textbf{O que isso significou:} o self-repair virou um grafo de estados explícito,
  reusando os mesmos prompts e o mesmo teto de 1 tentativa — mudança estrutural, não de
  comportamento. Confirmado com 5 testes automatizados e teste manual do dashboard real.

  \pause
  \begin{decisaobloco}{O que decidimos}
    Promovido para produção. Mantido o teto de 1 tentativa do original — o protótipo com
    teto de 6 mostrou o modelo local repetindo o mesmo erro sem se corrigir de fato.
  \end{decisaobloco}
\end{frame}

\begin{frame}{Roteador: Python puro vs. LangGraph}
  \small
  \begin{numerobloco}{O que medimos (golden set completo, 67 perguntas)}
    \begin{center}
    \begin{tabular}{lc}
      \toprule
      \textbf{Cenário} & \textbf{Divergências} \\
      \midrule
      Sem desempate por LLM (\texttt{usar\_llm=False}) & 0/67 \\
      Com desempate por LLM real (\texttt{usar\_llm=True}) & 0/67 \\
      \bottomrule
    \end{tabular}
    \end{center}
  \end{numerobloco}

  \pause
  \textbf{O que isso significou:} os 11 \emph{answerability gates} determinísticos e o
  desempate por LLM foram reusados sem reescrever nenhuma regra — só a orquestração virou
  um \texttt{StateGraph} de 2 nós (gates em cascata + desempate condicional) em vez da
  cadeia de \texttt{if} original.

  \pause
  \begin{decisaobloco}{O que decidimos}
    Promovido para produção. Item de maior risco/esforço do roadmap (9+ gates a portar sem
    perder cobertura) fechado sem nenhuma regressão medida.
  \end{decisaobloco}
\end{frame}

\begin{frame}{Achado ao testar em produção real: bugs de import}
  \begin{numerobloco}{O problema}
    \small
    \texttt{streamlit run app.py} (fora do root do projeto no \texttt{sys.path}) quebrou
    com \texttt{ModuleNotFoundError} em cadeia — nome de módulo ambíguo entre ``módulo
    solto'' e ``pacote'', dependendo de quem importa primeiro.
  \end{numerobloco}

  \pause
  \textbf{O que isso significou:} testes automatizados sempre rodaram com o root do
  projeto já no \texttt{sys.path} por acidente — escondendo o bug até o teste manual.

  \pause
  \begin{decisaobloco}{O que decidimos}
    \small
    Carregar os módulos migrados por \textbf{caminho de arquivo explícito}
    (\texttt{importlib.util.spec\_from\_file\_location}) — elimina a ambiguidade,
    confirmado nos dois cenários e no dashboard ao vivo.
  \end{decisaobloco}
\end{frame}

\begin{frame}{Próximos passos}
  \begin{itemize}
    \item Diagnóstico em camadas $\to$ avaliar se vale migrar (baixo retorno esperado).
    \item NL-to-SQL (geração) $\to$ maior risco, validação determinística própria precisa
      ser preservada.
    \item Agentic RAG (retrieval adaptativo) — protótipo medido, \alert{resultado negativo}
      (Recall@1 caiu de 12/15 para 11/15), não promovido; causa raiz: chunk certo às vezes
      não está no topo do documento certo, reformular a query não resolve isso.
    \item Contextual Retrieval — protótipo medido, resultado misto (MRR +0,008,
      Precision@5 $-3,6$pp), não promovido; manuais curtos não se beneficiam tanto quanto
      a literatura original (documentos longos) prevê.
    \item Regra que não muda: nenhuma migração é aceita se regredir
      \texttt{eval/rodar\_golden.py} (roteamento/faithfulness) ou o golden set.
  \end{itemize}
\end{frame}

\end{document}
